\section{后期拟完成的研究工作及进度安排}

\subsection{后期拟完成的工作}

中期阶段之后，本课题的后续工作将主要围绕实验补齐、训练推进、系统收口和论文定稿四个方面展开。首先，需要继续完成系统效果验证中的缺失部分，其中最核心的任务是补齐横向对比实验。当前系统已经完成了纵向实验框架、阶段性运行和治理复核，但横向对照尚未形成正式结果，因此后续必须在统一实验条件下完成与典型代理范式的对比运行，并对执行成功率、恢复行为、人工介入情况、审计完整性和端到端稳定性等指标进行整理。只有在这一部分补齐之后，论文中的实验章节才能从“阶段性说明”进一步推进到“具有完整比较依据的正式论证”。

其次，需要推进外部模型专用训练及其后续评估工作。当前系统已经完成训练数据抽取、质量门禁和训练方案准备，但外部模型训练本身尚未正式启动，因此后续需要在统一数据口径下完成训练实施、训练结果整理和训练后评估，并进一步验证模型输出是否能够在候选质量、可执行性和恢复稳定性方面形成可解释的提升。与此相关的另一项工作，是继续完善离线评估口径与统一门禁规则，使训练结果、系统运行结果和实验报告之间能够保持一致的评价标准，从而避免出现“模型已训练但无法形成稳定结论”的情况。

再次，在系统实现层面，后续仍需继续完成多工具统一验收、自动化门禁完善以及部分边界场景补样等工作。当前系统主链已经成型，但若要支撑最终论文定稿和答辩展示，还需要进一步提高工具接入的一致性、运行稳定性和证据完整性，尤其是在多工具场景下的统一调用规范、异常处理策略和验收标准方面，需要继续做收口工作。与此同时，已有的图示、日志证据和阶段性报告还需要进一步整理为可直接用于论文插图、实验图表和答辩展示的标准材料。

最后，论文写作本身也将进入由“章节骨架搭建”转向“正式内容收束”的阶段。除继续完善后续章节正文之外，还需要在实验结论、图表索引、语言风格统一和章节衔接方面做进一步整理，使论文从当前的中期报告版本逐步过渡到结构完整、论证闭环清晰的正式论文版本。换言之，后期工作的重点已经不再是扩展大量新的功能点，而是围绕已有系统成果，将实验、训练、评估和论文表达统一到同一条清晰主线之上。

\subsection{后期工作的进度安排}

从推进顺序上看，后续工作应当按照“先补实验与评估，再推进训练与收口，最后完成论文定稿与答辩准备”的节奏展开。第一阶段的重点是实验补齐与评估收口，即优先完成横向对比实验、统一已有纵向实验与治理复核结果，并同步整理阶段性图表和指标说明。这一阶段的目标不是继续扩展系统范围，而是尽快把当前已经具备的实验框架转化为可直接写入论文的正式结果，使中期之后的工作建立在更稳定的证据基础之上。

第二阶段的重点是完成外部模型专用训练及相关验证，并同步推进多工具统一验收和门禁完善。由于训练尚未启动，这部分工作必须在实验口径基本稳定之后进行，以避免训练输入、评估方式和系统结论之间出现脱节。在这一阶段中，需要一并完成训练结果复核、必要的回退策略验证，以及多工具场景下的运行一致性检查，从而保证后续论文中的“系统实现”“实验评估”和“局限性分析”三部分能够互相对应。

第三阶段的重点是论文与答辩材料的最终定稿。届时需要在已有章节基础上继续完善实验章节、问题分析、结论与展望等内容，并统一图示、表格、术语和章节语言，使全文形成风格一致、逻辑闭环清晰的正式稿。同时，还需要围绕系统总览图、核心机制图、实验表格和阶段性结论准备答辩展示材料，使论文文本与答辩表达能够保持一致。总体而言，后期工作的边界已经较为清楚，剩余任务虽然仍然较多，但主要属于在既有系统和既有框架基础上的补齐、收口与整合，因此具备继续推进并按时完成的现实基础。
